\documentclass[a4paper,10pt]{article}

\usepackage[utf8]{inputenc}
\usepackage[T1]{fontenc}
\usepackage[turkish]{babel}
\usepackage{geometry}
\usepackage{booktabs}
\usepackage{array}
\usepackage{tabularx}
\usepackage{longtable}
\usepackage{enumitem}
\usepackage{listings}
\usepackage{xcolor}
\usepackage[hidelinks]{hyperref}
\usepackage{microtype}
\usepackage{amsmath}
\usepackage{titlesec}

% Türkçe babel paketinin "=" karakterini listings seçeneklerinde etkilemesini engeller
\AtBeginDocument{\shorthandoff{=}}

\geometry{
    a4paper,
    top=1.8cm,
    bottom=1.8cm,
    left=1.8cm,
    right=1.8cm
}

\setlength{\parindent}{0pt}
\setlength{\parskip}{5pt}
\renewcommand{\arraystretch}{1.18}

\titleformat{\section}{\large\bfseries}{\thesection.}{0.5em}{}
\titleformat{\subsection}{\normalsize\bfseries}{\thesubsection.}{0.5em}{}

\lstdefinestyle{cstyle}{
    language={[ANSI]C},
    basicstyle=\ttfamily\small,
    keywordstyle=\bfseries,
    commentstyle=\itshape,
    numbers=left,
    numberstyle=\tiny,
    numbersep=7pt,
    frame=single,
    breaklines=true,
    showstringspaces=false,
    tabsize=4
}

\newcommand{\hex}[1]{\texttt{0x#1}}
\newcommand{\code}[1]{\texttt{#1}}

\title{\textbf{DEOS LoRa Telemetry ICD}\\
\large Sürüm 1.0}
\author{DEOS -- Dokuz Eylül Otonom Sistemler}
\date{\today}

\begin{document}

\maketitle

\section{Giriş ve Kapsam}
Bu doküman (ICD), DEOS araç içi ağı (CAN-FD / Ethernet) ile \textbf{Yer Kontrol İstasyonu} arasında kurulan \textbf{LoRa (Uzun Mesafe Radyo)} tabanlı telemetri haberleşmesinin kurallarını, çerçeve (framing) yapısını ve hata telafi (ACK/Retry) algoritmalarını tanımlar.

Yer Kontrol yazılımını geliştirecek mühendislerin, LoRa modülüne seri port (UART) üzerinden basacakları veriyi bu dokümanda belirtilen kurallara göre şekillendirmesi zorunludur.

\section{Temel Mimari ve Şeffaflık (Transparency)}
DEOS mimarisinde tüm mesajlar \code{deos\_message\_t} adlı 84 baytlık standart bir yapıya (struct) sahiptir.
LoRa arayüzü, bu standart DEOS yapısını \textbf{hiç bozmadan} (şeffaf bir şekilde) havadan taşımakla yükümlüdür. Yer Kontrol İstasyonu, kendi içinde bir \code{deos\_message\_t} paketi oluşturur, bunu bir ``LoRa Çerçevesi (Frame)'' içine hapseder ve gönderir. Araçtaki Main Node (Gateway), bu çerçeveyi çözer, içinden orijinal DEOS mesajını çıkarır ve doğrudan aracın iç ağına (Router) yönlendirir.

\section{LoRa Çerçeve (Frame) Yapısı}
Seri haberleşmede byte kaymalarını (misalignment) engellemek ve havadaki veri kayıplarını saptamak için her mesaj bir \textbf{Senkronizasyon Başlığı (Header)} ve \textbf{CRC16} doğrulaması ile paketlenir.

Genel paket yapısı şu şekildedir:
\begin{table}[h]
\centering
\begin{tabular}{|c|c|c|c|c|c|}
\hline
\textbf{MAGIC\_1} & \textbf{MAGIC\_2} & \textbf{Flag} & \textbf{Length} & \textbf{Payload (DEOS Mesajı)} & \textbf{CRC16} \\
\hline
1 Byte & 1 Byte & 1 Byte & 1 Byte & $N$ Byte & 2 Byte \\
\hline
\hex{AA} & \hex{55} & (Bkz. Bayraklar) & $N$ & Full DEOS Frame & CCITT \\
\hline
\end{tabular}
\end{table}

\subsection{Alanların Açıklamaları}
\begin{itemize}[label=\textbullet]
    \item \textbf{MAGIC\_1 \& MAGIC\_2:} Senkronizasyon baytlarıdır. Sabit olarak \hex{AA} ve \hex{55} değerlerini alırlar. Alıcı taraf, UART üzerinden gelen veri akışında bu iki baytı yan yana görene kadar tüm verileri çöpe atmalıdır.
    \item \textbf{Flag (Bayrak):} Mesajın türünü ve ACK talebini belirtir.
    \item \textbf{Length:} Payload alanının uzunluğunu (bayt cinsinden) belirtir. Dinamik olarak hesaplanır (\code{offsetof(deos\_message\_t, payload) + payload\_len}). ACK paketlerinde ise sadece CRC doğrulama verisi taşındığı için \textbf{2} olur.
    \item \textbf{Payload (Full DEOS Frame):} Esas verinin taşındığı yerdir. Bu alan \code{deos\_message\_t} yapısını taşır. Gönderim formatı şu şekildedir: Struct'ın başından \code{payload} dizisine kadar olan sabit kısım (Header alanı) kopyalanır, ardından sadece \code{payload\_len} kadar dolu olan gerçek veri baytları eklenir. Kalan boş baytlar hava süresi (airtime) tasarrufu için gönderilmez.
    \item \textbf{CRC16:} Sırasıyla Flag, Length ve Payload alanları üzerinden hesaplanan 16-bit CRC-CCITT (Polinom: 0x1021) değeridir. \textit{MAGIC baytları CRC hesaplamasına dahil edilmez.} (Little-endian olarak, önce LSB sonra MSB yollanır).
\end{itemize}

\subsection{Flag (Bayrak) Değerleri}
\begin{table}[h]
\centering
\begin{tabularx}{\textwidth}{|l|c|X|}
\hline
\textbf{Makro Adı} & \textbf{Değer} & \textbf{Açıklama} \\
\hline
\code{LORA\_FLAG\_NONE} & \hex{00} & Normal bir veri paketidir. Alıcı bu paketi onaylamak (ACK göndermek) zorunda değildir. \\
\hline
\code{LORA\_FLAG\_ACK\_REQ} & \hex{01} & Alıcıdan bu paket için kesinlikle bir Alındı Onayı (ACK) beklenmektedir. \\
\hline
\code{LORA\_FLAG\_IS\_ACK} & \hex{02} & Bu paket bir onay (ACK) paketidir. Payload kısmında, onaylanan paketin CRC16 değeri (2 byte) yer alır. \\
\hline
\end{tabularx}
\end{table}

\section{ACK ve Tekrar Gönderme (Retry) Mekanizması}
Araçtan yere veya yerden araca giden bazı hayati mesajlar (örneğin Acil Durum / E-Stop, Yörünge Gönderme vb.) havada kaybolmamalıdır. Bu tür kritik iletimlerde Yer Kontrol yazılımı şu senaryoyu işletmelidir:

\subsection{Gönderici (Yer Kontrol) Davranışı}
\begin{enumerate}
    \item Yer kontrol yazılımı \code{deos\_message\_t} paketini hazırlar.
    \item Paketi LoRa çerçevesine sokarken Flag alanını \hex{01} (\code{LORA\_FLAG\_ACK\_REQ}) olarak ayarlar ve CRC'sini hesaplayarak UART'tan basar.
    \item Gönderim tamamlandığı an bir zamanlayıcı (timer) başlatılır. Maksimum bekleme süresi (Timeout) genellikle 500ms ile 1000ms arası seçilmelidir.
    \item Bu süre içinde alıcıdan \code{LORA\_FLAG\_IS\_ACK} bayrağına sahip bir onay paketi gelmezse, veya gelen onay paketinin içindeki (Payload) CRC değeri, bizim gönderdiğimiz paketin CRC'si ile eşleşmezse; paket \textbf{kayıp} sayılır.
    \item Kayıp durumunda aynı çerçeve tekrar yollanır. Bu işlem belirlenen maksimum tekrar sayısına (Örn: 3 kez) kadar denenir.
\end{enumerate}

\subsection{Alıcı (Araç - Main Node) Davranışı}
\begin{enumerate}
    \item Araçtaki sistem, UART üzerinden gelen baytları dinler. \hex{AA 55} sekansını yakaladığında okumaya başlar.
    \item Gelen Flag değerini okur, ardından Length değerine göre Payload'ı okur, sonunda CRC16 değerini kendi hesapladığı CRC ile karşılaştırır. Eşleşmezse paketi yok sayar.
    \item Eğer eşleşirse ve Flag değeri \hex{01} (\code{ACK\_REQ}) ise; araç hiç vakit kaybetmeden (DEOS Router'a mesajı iletmeden önce) bir ACK paketi hazırlar.
    \item ACK Paketinin yapısı: \\ \code{[0xAA] [0x55] [Flag: 0x02] [Length: 0x02] [Gelen Mesajın CRC'si (2 Byte)] [Yeni CRC16 (2 Byte)]} \\ şeklindedir ve derhal Yer Kontrol'e geri yollanır.
\end{enumerate}

\section{Örnek CRC-CCITT (0x1021) C Algoritması}
Yer Kontrol yazılımını geliştirecek taraf (Python, C++, C\# vb.), paket bütünlüğünü kontrol etmek için standart CCITT algoritmasını kullanmalıdır. C tabanlı bir örnek aşağıda verilmiştir:

\begin{lstlisting}[style=cstyle]
uint16_t crc16_ccitt(uint16_t crc, const uint8_t *src, size_t len)
{
    while (len--) {
        crc = crc ^ (*src++ << 8);
        for (int i = 0; i < 8; i++) {
            if (crc & 0x8000)
                crc = (crc << 1) ^ 0x1021;
            else
                crc = crc << 1;
        }
    }
    return crc;
}
\end{lstlisting}
\textit{Not: Yukarıdaki fonksiyona başlangıç değeri (seed) olarak \hex{FFFF} verilmelidir. Hesaplamaya \code{Flag} baytından başlanmalı, ardından \code{Length} ve Payload baytları ile devam edilmelidir. MAGIC baytları atlanmalıdır.}

\end{document}
